\section{Проектирование гибридной интеллектуальной системы}

\subsection{Определение пользователей гибридной интеллектуальной системы}
% Определяются категории пользователей и внешних подсистем, взаимодействующих
% с гибридной интеллектуальной системой.
%
% Для каждой категории указываются роль, цели, уровень подготовки,
% полномочия, доступные функции и ожидаемые результаты взаимодействия.
%
% Например: эксперт предметной области, оператор, конечный пользователь,
% администратор базы знаний, администратор системы, внешний программный
% сервис, внешний источник данных.

\subsection{Сценарии использования гибридной интеллектуальной системы}
% Описываются сценарии использования (use cases), отражающие взаимодействие
% пользователей и внешних систем с гибридной интеллектуальной системой.
%
% Для каждого сценария рекомендуется указать:
% - инициатора;
% - предусловия;
% - входные данные;
% - основной поток действий;
% - взаимодействие базы знаний, решателя задач, интерфейса и компонентов
%   логико-символьного и нейробионического подходов;
% - альтернативные, граничные и ошибочные потоки;
% - итоговый результат.
%
% Следует выделить типичные, граничные и ошибочные сценарии, а также
% сквозные сценарии от пользовательского запроса до объяснимого результата.

\subsection{Общая архитектура гибридной интеллектуальной системы}
% Представляется общая архитектура гибридной интеллектуальной системы
% и её место в контексте предметной области.
%
% Рекомендуется привести компонентную или структурную диаграмму,
% показывающую базу знаний, решатель задач, пользовательский интерфейс,
% компонент нейробионического подхода, внешние сервисы, хранилища данных
% и направления обмена данными.
%
% Следует показать, как архитектура обеспечивает выполнение сквозных
% сценариев и независимое обновление компонентов.

\subsection{Проектирование алгоритмов и методов решения задач}
% Детально описываются методы и алгоритмы, используемые в компонентах
% гибридной интеллектуальной системы.
%
% Для логико-символьного подхода могут описываться логический вывод,
% продукционный вывод, поиск, проверка ограничений, планирование,
% формирование объяснений.
%
% Для нейробионического подхода могут описываться классификация,
% ранжирование, извлечение признаков, поиск по эмбеддингам, генерация,
% обработка естественного языка или иной выбранный метод.
%
% Указываются входные и выходные данные, этапы алгоритмов, условия
% переходов, критерии остановки, порядок обращения к базе знаний,
% правила интеграции результатов и порядок обработки ошибок.
%
% Алгоритмы рекомендуется представить в виде псевдокода, блок-схем,
% диаграмм деятельности или диаграмм последовательности.
% Анализируется вычислительная сложность проектируемых алгоритмов
% и компонентов гибридной интеллектуальной системы.
%
% Рассматривается влияние объёма базы знаний, количества правил,
% глубины поиска, числа ограничений, объёма данных, размерности эмбеддингов,
% сложности нейробионической модели и числа параллельных запросов на
% время выполнения и потребление вычислительных ресурсов.
%
% Формулируются проектные ограничения и меры по их снижению:
% кэширование, индексация знаний, предварительная фильтрация,
% ограничение глубины поиска, эвристики, использование векторных индексов,
% ограничение длины контекста, пакетная обработка, декомпозиция задачи
% или асинхронное выполнение компонентов.


\subsection{Выводы}
% Кратко фиксируются результаты проектирования: пользователи и сценарии,
% архитектура, стратегия интеграции подходов, концептуальная и логическая
% модели, алгоритмы, форматы данных и проектные ограничения.
